%% Need to keep this so stuff works
\DocumentMetadata{tagging=on,
    pdfversion = 2.0,
    tagging-setup={math/setup=mathml-SE},
    pdfstandard=ua-2,
    lang=en-US
}

\documentclass[12pt, oneside, letterpaper]{book}
%%
%
%%
%%% DOCUMENT START
\input{preamble} % Moved to another file for ease of updating
\begin{document}
%%
%
%%
%%% Document Setup Settings
%%
\pagestyle{fancy} % Set the page style to "fancy"...
\renewcommand{\headrulewidth}{0pt}
\fancyhead{}
\fancyhead[R]{\thepage}
\fancyfoot{} % clear all footer fields
\renewcommand{\footrulewidth}{0.1pt}
\fancyfoot[L]{Truman Science Center}
\fancyfoot[C]{Calorimetry/Heat Capacity}
\fancyfoot[R]{15 April 2026}
%%
%
%% Title Set up
\begin{flushleft}
	\vspace{.3cm}
	{\textbf { \Large Calorimetry/Heat Capacity}}
\end{flushleft}
\vspace{1mm}
	\hrule
%%
%
%%
%%% START TO DOCUMENT TEXT
%%
%
The equation below should come to mind whenever you see the question mention calorimetry, heat, or specific heat. 
\begin{gather*}
    Q=mc\ \Delta T
\end{gather*}

\section{The Process}

\begin{enumerate}
   \item Assign any values given in the question to variables in the heat equation.
   \item Check the Units!!
   \begin{enumerate}
     \item Make sure the values you have for heat, mass, and temperature are the same as the units given in the specific heat. 
     \item Make sure your final answer is in the same units as the question asked for.
    \end{enumerate}
  \item Rearrange the heat equation so the variable you want to learn is isolated. 
  \item Plug in values 
\end{enumerate}

\subsection{Helpful Reference Values}

\subsubsection{Values for Water}
\begin{align*}
    &   c_{water}\ = 4.184\ \cfrac{\si{\J}}{\si{\g\degreeCelsius}} & d_{water}\ = 1\ \cfrac{\si{\g}}{\si{\mL}} &
\end{align*}

\subsubsection{Unit Conversions}
\begin{align*}
    & \SI{1}{\Cal} = \SI{1000}{\cal} & \SI{1}{\cal} = \SI{4.184}{\J} && \si{\degreeCelsius}\ = (\ \si{\degree} F - 32)\  \cdot \cfrac{5}{9} & 
\end{align*}

\newpage
\parskip = 6cm

\section{Calorimetry Practice} 

1) A $\qty{97.00}{\degreeCelsius}$ glass of water at $\qty{12.00}{\degreeCelsius}$ is left on the countertop. The next day the temperature of the water is $\qty{23.00}{\degreeCelsius}$, how much heat (in Joules) was absorbed by the water?


2) You have a special microwave button that releases $\qty{7.78}{\kilo\joule}$ of heat into your food. You feel like playing Hot Potato so you put a $\qty{25}{\degreeCelsius}$ potato into the microwave and turn it on. When it comes out it is $\qty{67}{\degreeCelsius}$ and it burns your hand. If the specific heat capacity of potatoes is $\qty{3.43}{  \cfrac{\joule}{\gram\degreeCelsius}}$, what is the mass of the potato?


3) You left a $\qty{3.11}{\gram}$ pure copper penny in your car which the sun heated to $\qty{32.00}{\degreeCelsius}$. When you took the penny out it released $\qty{10.78}{J}$ of heat and its temperature went down to $\qty{23.00}{\degreeCelsius}$. What is the specific heat capacity of Copper?


4) You pick up a $\qty{0.24}{\kilo\gram}$ rock from a cold, $50 \Fahrenheit$ natural spring. As it sits in the sun it heats up to $73.4 \Fahrenheit$. It's specific heat capacity is $\qty[per-mode=fraction]{910}{\milli\joule\per\gram\per \degreeCelsius}$, how much heat (in $\si{\kilo\joule}$) did the rock absorb?


5) You heat up a cast iron skillet that has a mass of $\qty{2.23}{\kilo\gram}$ to a temperature of $\qty{218}{\degreeCelsius}$. How many $\si{\kilo\joule}$ of energy are released from the skillet when it cools to room temperature ($\qty{22}{\degreeCelsius})$? The specific heat capacity of cast iron is $\qty[per-mode=fraction]{0.461}{\joule\per\gram\per\degreeCelsius}$


6) It takes $\qty{874.6}{\kilo\joule}$ to heat up $\qty{3.12}{\kilo\gram}$ of water to its boiling point at $\qty{100}{\degreeCelsius}$. What was the initial temperature of the water?

\newpage

\section{More Challenging Calorimetry Problems}
7) A metal object releases $\qty{67.78}{\kilo\joule}$ of energy into $\qty{300} {\mL}$ of water at $\qty{23} {\degreeCelsius}$. What temperature does the water rise to assuming the density of water is $\qty[per-mode=fraction]{1}{\gram\per{\milli\liter}}$

8) When you set a Dorito on fire over a can of $\qty{50.0} {\gram}$ of water, it releases $\qty{12.72}{\Cal}$ of energy. Assuming all of the energy goes into the water (and none into the can) what is the rise in temperature of the water?

9) a $\qty{20.0}{\gram}$ sample of a metal alloy at $\qty{82.0}{\degreeCelsius}$ is placed into $\qty{80.0}{\gram}$ of water at $\qty{25.3}{\degreeCelsius}$. The final temperature of the water is $\qty{29.3}{\degreeCelsius}$. Assuming there is no loss of heat to the surroundings, calculate the heat capacity of the alloy.

10) $\qty{250}{\mL}$ of water at $\qty{93}{\degreeCelsius}$ is mixed with $\qty{100}{\mL}$ fo water at $\qty{23} {\degreeCelsius}$. Calculate the final temperature of this mixture.

11) $\qty{30.5}{\gram}$ of metal at $\qty{105.0}{\degreeCelsius}$ is placed into $\qty{200}{\mL}$ of water at $\qty{24.00} {\degreeCelsius}$. If the final temperature of the water is $\qty{25.1}{\degreeCelsius}$, calculate the specific heat of the metal and identify it.

12) An unknown mass of iron at $\qty{100}{\degreeCelsius}$ is placed in an isolated cup containing $\qty{75.0}{\mL}$ of water at $\qty{40.0}{\degreeCelsius}$. The final temperature of the iron bolt is $\qty{41.3}{\degreeCelsius}$. Assuming that the heat capacity of the container is negligible, calculate the mass of the iron bolt.

\end{document}